% \documentclass{report}

% \usepackage{graphicx} \usepackage{dcolumn} \usepackage{bm}
% \usepackage{amsfonts, amssymb, amsmath} \usepackage{braket}

%\newcommand{\ket}[1]{\textstyle \left| #1 \right\rangle \displaystyle}
%\newcommand{\bra}[1]{\textstyle \left\langle #1 \right| \displaystyle}
%\newcommand{\braket}[2]{\textstyle \left\langle #1 \left|  #2
%\right\rangle\right. \displaystyle} \newcommand{\im}{\mathrm{i}}
%\newcommand{\e}{\e} \newcommand{\pwrt}[2]{\frac{\partial #1}{\partial #2}}
%\newcommand{\diff}{\mathrm{d}}
%
%\begin{document}

\subsection{Symmetry concerns within the silicon wavefunctions.}

We expect our full wavefunction to have the form
\begin{equation}
\Psi^{\left(\kappa,i\right)}
  =
    \sum_{p}
    \Phi_{p}\mathrm{e}^{\im\bm{k}_{p}\cdot\bm{r}}u_{\bm{k}_{p}\bm{r}}^{}
\label{form:full-symmetric-wavefunction}
\end{equation}
The transformation of the Bloch function terms above is straight-forward, but
the symmetry of the full wavefunction involves symmetrization of the product of
the Bloch and envelope terms. To understand the symmetry properties of the
envelope function requires us to understand the symmetrization of the full
wavefunction.

To understand the symmetrization of our full wavefunction, we apply our
projection symmetrization operator defined in \eqref{symm-operator} to a general
wavefunction. We assert that the full wavefunction is already symmetrized, and
as such, the operator leaves our full wavefunction invariant. In essence, we are
asserting that $\Psi$ is such that
\begin{equation}
  \Psi = \mathcal{P}_{ii}^{\left(\kappa\right)}\Psi.
\end{equation}
By expanding $\Psi$ in terms of the form articulated above in
\eqref{form:full-symmetric-wavefunction} and matching terms we can assert
symmetrization conditions upon the envelope functions.

Our first step involves the following form, namely
\begin{equation}
\mathcal{V}_{\kappa,i}^{\dagger}\Psi
  =
    \left(\alpha_{0}^{\left(\kappa,i\right)}\right)^{*}
    \left(xyz\right)
    \Psi
    +
    \left(\alpha_{1}^{\left(\kappa,i\right)}\right)^{*}
    \left(zxy\right)
    \Psi
    +
    \left(\alpha_{2}^{\left(\kappa,i\right)}\right)^{*}
    \left(yzx\right)
    \Psi.
\end{equation}
It is useful to examine the operators included in $ \mathcal{V} _{\kappa,i}
^{\dagger} $ above. To wit, we can define a set of ancillary functions based off
an initial function. We describe these functions as being "valley indexed" by a
particular crystallographic axis. 
\begin{align}
  f_{z} & \equiv \left(xyz\right)f, \label{idx:z}\\
  f_{x} & \equiv \left(yzx\right)f, \label{idx:x}\\
  f_{y} & \equiv \left(zxy\right)f. \label{idx:y}
\end{align}
It is useful to note these ancillary functions have the following relationships,
\begin{align}
  f_{z} & = \left(zxy\right)f_{x} = \left(yzx\right)f_{y}, \\
  f_{x} & = \left(yzx\right)f_{z} = \left(zxy\right)f_{y}, \\
  f_{y} & = \left(zxy\right)f_{z} = \left(yzx\right)f_{x}.
\end{align}
For ease of nomenclature, we'll call these three transformations
$\left(xyz\right)$, $\left(yzx\right)$, and $\left(zxy\right)$ the indexing
transformations.

Since we've introduced the concept of valley indexed functions, we'll introduce
an additional stipulation upon $\Psi$, namely that
\begin{equation}
  \Psi = c_{z}\psi_{z} + c_{x}\psi_{x} + c_{y}\psi_{y}.
\label{wf_raw}
\end{equation}
Above, the three functions $\psi_{x}$, $\psi_{y}$, and $\psi_{z}$ are arbitrary
but linked by valley indexing relationships introduced in \eqref{idx:z},
\eqref{idx:x}, and \eqref{idx:y}. To be consistent with definition
\eqref{wf_raw} and equation \eqref{form:full-symmetric-wavefunction} we expect
the valley-linked $\psi_{q}$ terms of the full wavefunction to themselves be
constructed from terms involving the product of the envelope and valley-indexed
Bloch functions. We can group those terms by Bloch function axes, and so we
expect the following
\begin{equation}
  \psi_{z} = \psi_{\parallel,z} + \psi_{\perp_{1},x} + \psi_{\perp_{2},y},
\label{single-term:form}
\end{equation}
where 
\begin{equation}
  \psi_{\aleph,q}
    =
    \Phi_{\aleph,+q}
    \theta^{\left(+q\right)}
      +
    \Phi_{\aleph,-q}
    \theta^{\left(-q\right)}
.
\label{single-axis:form0}
\end{equation}
where the Bloch functions $\theta^{\left(q\right)}$ are given as
\begin{equation}
  \theta^{\left(q\right)} 
  =
    \mathrm{e}^{\im\bm{k}_{q}\cdot\bm{r}}
    u_{\bm{k}_{q}\bm{r}}^{}
.
\end{equation}

The indexing relations may be used to construct $\psi_{x}$ and $\psi_{y}$ from
$\psi_{z}$. We may alternatively write the valley-indexed terms of the full
wavefunction as
\begin{equation}
  \psi_{q}
  =
  \psi_{\aleph_{qz},z} + \psi_{\aleph_{qx},x} + \psi_{\aleph_{qy},y}.
\tag{\ref{single-term:form}a}
\label{single-term:form-a}
\end{equation}
with possible values for $\aleph_{qq^{\prime}}$ going as $\aleph_{qq} =
\parallel$, $\aleph_{qr} = \perp_{1}$ for $qr \in \left\{yz,zx,xy\right\}$, and
$\aleph_{qs} = \perp_{2}$ for $qs \in \left\{xz,yx,zy\right\}$.

When we apply the first factor operator of the symmetrization projection
operator to the form in \eqref{wf_raw}, we'll find different results depending
on the particular symmetry we are projecting upon. For $\kappa = A,E$ we can
note the following relation involving the $\alpha$ coefficients, with
\begin{equation}
9\alpha_{1}\alpha_{2} = 1.
\end{equation}
If, in addition to the relationship above, we also note the following relation
\begin{equation}
\mathcal{V}_{\kappa,i}^{\dagger}\psi_{z}
  =
3\alpha_{1}^{\left(\kappa,i\right)}\mathcal{V}_{\kappa,i}^{\dagger}\psi_{x}
  =
3\alpha_{2}^{\left(\kappa,i\right)}\mathcal{V}_{\kappa,i}^{\dagger}\psi_{y}
,
\end{equation}
we may find for a $\kappa = A,E$ state that
\begin{equation}
\mathcal{V}_{\kappa,i}^{\dagger}\Psi
  =
  \left(c_{z}+3c_{x}\alpha_{2}+3c_{y}\alpha_{1}\right)
  \left(
    \alpha_{0}^{\left(\kappa,i\right)}\psi_{z}
    +
    \alpha_{1}^{\left(\kappa,i\right)}\psi_{x}
    +
    \alpha_{2}^{\left(\kappa,i\right)}\psi_{y}
  \right)
  .
\end{equation}
This shows a clear projection form, with a resulting symmetrized function form
multiplied by a overall factor indicating how much of the original function can
be said to "lie along" the symmetrized form. For $\kappa = A,E$ then, we can
take the $c$ coefficients to have the following form, with $ c_{z} = C/3 $, $
c_{x} = C \alpha_{ 1 } $, and $c_{y} = C \alpha_{ 2 } $, where $C$ is a complex
factor useful for normalization. If we define the $c_{q}$ coefficients as such,
we'll find that
\begin{equation}
\mathcal{V}_{\kappa,i}^{\dagger}
\Psi
  =
\Psi
  =
C
\left(
  \alpha_{0}^{\left(\kappa,i\right)}\psi_{z}
  +
  \alpha_{1}^{\left(\kappa,i\right)}\psi_{x}
  +
  \alpha_{2}^{\left(\kappa,i\right)}\psi_{y}
\right)
\label{wf_symm:a}
\end{equation}

In the case of $\kappa = T$, the $\alpha$ go as
\begin{gather}
  \alpha_{0}^{\left(T_{n},i\right)} = \delta_{iz}, \quad\quad
  \alpha_{1}^{\left(T_{n},i\right)} = \delta_{ix}, \quad\quad
  \alpha_{2}^{\left(T_{n},i\right)} = \delta_{iy}.
  % \alpha_{q}^{2} &= \alpha_{q}, \\
  % \alpha_{q}^{*} &= \alpha_{q}, \\
  % \alpha_{q}\alpha_{r} &= \delta_{iq}\delta_{qr}.
\end{gather}
This means application of the $\mathcal{V}^{\dagger}$ operator yields the
following effect upon our wavefunction, namely
\begin{align}
  \mathcal{V}_{T_{n},i}^{\dagger}
  \Psi
  & =
  d_{z}\psi_{z}
  +
  d_{x}\psi_{x}
  +
  d_{y}\psi_{y}
  ,
\label{wf_symm:b}
\\
  d_{z}
  & =
  \delta_{iz}c_{z}
  +
  \delta_{ix}c_{x}
  +
  \delta_{iy}c_{y}
  ,
\\
  d_{x}
  & =
  \delta_{iz}c_{x}
  +
  \delta_{ix}c_{y}
  +
  \delta_{iy}c_{z}
  ,
\\
  d_{y}
  & =
  \delta_{iz}c_{y}
  +
  \delta_{ix}c_{z}
  +
  \delta_{iy}c_{x}
  .
\end{align}
Alternatively, we can state
\begin{align}
  \mathcal{V}_{T_{n},i}^{\dagger}
  \Psi
  & =
  \delta_{iz}\psi_{0}
  +
  \delta_{ix}\psi_{1}
  +
  \delta_{iy}
  \psi_{2}
  ,
\label{wf_symm2}
\\
  \psi_{0} & = c_{z}\psi_{z}+c_{x}\psi_{x}+c_{y}\psi_{y},
\\
  \psi_{1} & = c_{x}\psi_{z}+c_{y}\psi_{x}+c_{z}\psi_{y},
\\
  \psi_{2} & = c_{y}\psi_{z}+c_{z}\psi_{x}+c_{x}\psi_{y}
  .
\end{align}
We see a marked difference between the effect of the
$\mathcal{V}_{T_{n},i}^{\dagger}$ factor operator and the
$\mathcal{V}_{A_{n},i}^{\dagger}$, $\mathcal{V}_{E,i}^{\dagger}$ operators. For
$\kappa = T$, in essence the effect of the valley mixing/rotation factor
operator $\mathcal{V}_{T_{n},i}^{\dagger}$ is to generate a cyclic permutation
of the $c$ coefficients without an overall change of function form. Another way
to interpret the effect of the $\mathcal{V}_{T_{n},i}^{\dagger}$ factor operator
is that it aligns the $z$-axis of the function with the $T$-states' axes of
symmetry. In the case of $\kappa = A,E$ we find that the factor operator
redistributes the weight of each valley.

With appropriate choice of the $\beta$ coefficients below, we can group results for all symmetries in the following form,
\begin{equation}
  \mathcal{V}_{\kappa,i}^{\dagger}
  \Psi
  =
  \beta_{z}^{\left(\kappa,i\right)}\psi_{z}
  +
  \beta_{x}^{\left(\kappa,i\right)}\psi_{x}
  +
  \beta_{y}^{\left(\kappa,i\right)}\psi_{y}.
\label{multi-term_form}
\end{equation}
We can group the terms solely by Bloch axes, which will give us 
\begin{align}
  \mathcal{V}_{\kappa,i}^{\dagger}
  \Psi
  &=
  \psi_{z}^{\prime} + \psi_{x}^{\prime} + \psi_{y}^{\prime},
\label{single-term:form-b}
\\
  \psi_{z}^{\prime}
  &=
  \beta_{z}^{\left(\kappa,i\right)}\psi_{\parallel,z}
  +
  \beta_{y}^{\left(\kappa,i\right)}\psi_{\perp_{1},z}
  +
  \beta_{x}^{\left(\kappa,i\right)}\psi_{\perp_{2},z}
  ,
\label{single-axis:form:z}
\\
  \psi_{x}^{\prime}
  &=
  \beta_{x}^{\left(\kappa,i\right)}\psi_{\parallel,x}
  +
  \beta_{z}^{\left(\kappa,i\right)}\psi_{\perp_{1},x}
  +
  \beta_{y}^{\left(\kappa,i\right)}\psi_{\perp_{2},x}
  ,
\label{single-axis:form:x}
\\
  \psi_{y}^{\prime}
  &=
  \beta_{y}^{\left(\kappa,i\right)}\psi_{\parallel,y}
  +
  \beta_{x}^{\left(\kappa,i\right)}\psi_{\perp_{1},y}
  +
  \beta_{z}^{\left(\kappa,i\right)}\psi_{\perp_{2},y}
  .
\label{single-axis:form:y}
\end{align}
We can rewrite this expression to make the underlying envelope-Bloch product
form explicit.
\begin{align}
\psi_{q}^{\prime}
  &=
    \sum_{\upsilon = +,-}
    \Phi_{\upsilon q}
    \theta^{\left(\upsilon q\right)}
,
\label{single-axis:form1}
\\
  \Phi_{\upsilon z}
  & =
  \beta_{z}^{\left(\kappa,i\right)}\Phi_{\parallel,\upsilon z}
  +
  \beta_{y}^{\left(\kappa,i\right)}\Phi_{\perp_{1},\upsilon z}
  +
  \beta_{x}^{\left(\kappa,i\right)}\Phi_{\perp_{2},\upsilon z},
\\
  \Phi_{\upsilon x}
  & =
  \beta_{x}^{\left(\kappa,i\right)}\Phi_{\parallel,\upsilon x}
  +
  \beta_{z}^{\left(\kappa,i\right)}\Phi_{\perp_{1},\upsilon x}
  +
  \beta_{y}^{\left(\kappa,i\right)}\Phi_{\perp_{2},\upsilon x}
\\
  \Phi_{\upsilon y}
  & =
  \beta_{y}^{\left(\kappa,i\right)}\Phi_{\parallel,\upsilon y}
  +
  \beta_{x}^{\left(\kappa,i\right)}\Phi_{\perp_{1},\upsilon y}
  +
  \beta_{z}^{\left(\kappa,i\right)}\Phi_{\perp_{2},\upsilon y}
\end{align}
It's worth noting that the symmetrizing operators can be extended down through
the layers of definition introduced in \eqref{single-term:form} and
\eqref{single-axis:form0} to the envelope functions $\Phi_{\aleph,\pm q}$
because the operators are linear. We see the $\beta$ coefficients show up in our
expressions above. The value of the $\beta$ coefficients depend upon the value
of the symmetry index $\kappa$ and follow directly from \eqref{wf_symm:a} and
\eqref{wf_symm:b}. There's an important distinction between $\beta$ for $\kappa
= T$ and $\beta$ for $\kappa = A,E$. For the $A,E$ states the coefficients are
not arbitrary, with $\beta_{q} = \alpha_{q}$. For the $T$ states, the value of
$\beta_{q} = d_{q}$ and thus depends upon the value of the $c$ arbitrary
constants. Since the factor operator for $\kappa = T$ does not remix the
valleys, the weight of each valley remains arbitrary. As such we note that the
envelope functions for states with $\kappa = T$ are not linked by the indexing
expressions introduced in \eqref{idx:z}, \eqref{idx:x}, and \eqref{idx:y}. This
is not the case for $\kappa = A,E$. Namely, we find the following
valley-indexing relationships.
\begin{align}
  \Phi_{\upsilon x}
  &=
  3\alpha_{1}^{\left(\kappa,i\right)}\left(yzx\right)
  \Phi_{\upsilon z},
\label{idx:env:xz}
\\
  \Phi_{\upsilon y}
  &=
  3\alpha_{2}^{\left(\kappa,i\right)}\left(zxy\right)
  \Phi_{\upsilon z},
\label{idx:env:yz}
\\
  \Phi_{\upsilon y}
  &=
  3\alpha_{1}^{\left(\kappa,i\right)}\left(yzx\right)
  \Phi_{\upsilon x}.
\label{idx:env:yx}
\end{align}

Over the course of this section, we will determine the necessary properties of
the envelope function to generate the desired symmetry. As we apply the rest of
the projection operator factors to our full wavefunction, we'll explore how the
interplay of the Bloch and envelope functions influences the process of
symmetrization.

Our next step will be to apply $\mathcal{A}_{z}^{\left(\kappa\right)}$ to our
wavefunction. We note that for $\kappa = E$, the operator
$\mathcal{A}_{z}^{\left(E\right)} = \left(xyz\right)$ and so the effect of the
operator is trivial. For $\kappa = A,T$, the effect of this operator upon our
terms are given below.

Without much trouble we can assert that
\begin{equation}
  \mathcal{A}_{z}^{\left(\kappa\right)}\psi_{z}^{\prime} = \psi_{z}^{\prime}.
\end{equation}
Given the linearity of the symmetry transformations, we can achieve this on a
per Bloch function term basis. For the $z$ term it is true that
\begin{equation} 
  \mathcal{A}_{z}^{\left(\kappa\right)}\Phi_{\upsilon z} = \Phi_{\upsilon z}
\label{acyc-z:env-z}
,
\end{equation}
if we also assert of our envelope function terms that
\begin{equation}
\Phi_{\upsilon z}
  =
    \sigma^{\left(\kappa\right)}\left(yxz\right)\Phi_{\upsilon z}
\label{acyc:z}
.
\end{equation}
However, there are subtleties when applying the acyclic operator
$\mathcal{A}_{z}$ to states with Bloch function factors indexed by $x$ and $y$.
Let's consider the following.
\begin{align}
  \mathcal{A}_{z}^{\left(\kappa\right)}
  \Phi_{\upsilon x}\theta^{\left(\upsilon x\right)}
  & =
  \frac
  {
    \Phi_{\upsilon x}
    \theta^{\left(\upsilon x\right)}
    + 
    \left[\sigma^{\left(\kappa\right)}\left(yxz\right)\Phi_{\upsilon x}\right]
    \theta^{\left(\upsilon y\right)}
  }
  {2}
\label{acyc-z:x}
,\\
  \mathcal{A}_{z}^{\left(\kappa\right)}
  \Phi_{\upsilon y}\theta^{\left(\upsilon y\right)}
  & =
  \frac
  {
    \Phi_{\upsilon y}
    \theta^{\left(\upsilon y\right)}
    + 
    \left[\sigma^{\left(\kappa\right)}\left(yxz\right)\Phi_{\upsilon y}\right]
    \theta^{\left(\upsilon x\right)}
  }
  {2}
\label{acyc-z:y}
\end{align}
Here we note that for $\kappa = T$ the envelope functions are not linked by the
indexing relationships. However, we assert that
\begin{equation}
  \mathcal{A}_{z}^{\left(T\right)}
  \left\{ 
    \Phi_{\upsilon x}\theta^{\left(\upsilon x\right)}
      +
    \Phi_{\upsilon y}\theta^{\left(\upsilon y\right)}
  \right\}
= 
  \Phi_{\upsilon x}\theta^{\left(\upsilon x\right)}
    +
  \Phi_{\upsilon y}\theta^{\left(\upsilon y\right)}
.
\label{acyc-z:env-xy}
\end{equation}
Combining \eqref{acyc-z:x} and \eqref{acyc-z:y} with \eqref{acyc-z:env-xy}
yields
\begin{align}
  \Phi_{\upsilon x}
  &=
  \sigma^{\left(\kappa\right)}\left(yxz\right)\Phi_{\upsilon y}
\label{acyc:x-y}
\\
  \Phi_{\upsilon y}
  &=
  \sigma^{\left(\kappa\right)}\left(yxz\right)\Phi_{\upsilon x}
\tag{\ref{acyc:x-y}${}^{\prime}$}
\label{acyc:y-x}
\end{align}
It should be noted that \eqref{acyc:x-y} and \eqref{acyc:y-x} are equivalent. We
can clearly state that for $\kappa = T$ states there exists a relationship
between the $\Phi_{\upsilon x}$ and $\Phi_{\upsilon y}$ envelopes, but no
relationship between these envelopes and the $\Phi_{\upsilon z}$ envelopes can be
derived.

However, we saw above in the indexing relations \eqref{idx:env:xz},
\eqref{idx:env:yz}, and \eqref{idx:env:yx} that relationships exist between the
envelope functions of different valley-indices for $\kappa = A,E$. We also
observe that
\begin{equation}
  \left(yzx\right)\left(yxz\right)\left(yzx\right)
=
  \left(zxy\right)\left(yxz\right)\left(zxy\right)
=
  \left(yxz\right).
\end{equation}
In particular, by combining the indexing relations with the observation above
and utilizing \eqref{acyc:z} it may be shown that
\begin{equation}
\left(yxz\right)\Phi_{\upsilon x}
  =
    3\alpha_{2}^{\left(\kappa,i\right)}\sigma^{\left(\kappa\right)}
    \Phi_{\upsilon y}
\label{acyc:x}
.
\end{equation}
or equivalently 
\begin{equation}
\left(yxz\right)\Phi_{\upsilon y}
  =
    3\alpha_{1}^{\left(\kappa,i\right)}\sigma^{\left(\kappa\right)}
    \Phi_{\upsilon x}
\tag{\ref{acyc:x}${}^{\prime}$}
\label{acyc:y}
.
\end{equation}
Since the $E$-states have a trivial acyclic operator, we'll keep our focus on
the $A$-states. In light of this consideration, we find $
3\alpha_{1}^{\left(A\right)} = 3\alpha_{2}^{\left(A\right)} = 1 $. If this is
true than \eqref{acyc:x-y} is equivalent to \eqref{acyc:x} and \eqref{acyc:y-x}
is equivalent to \eqref{acyc:y}. As such, it is fairly trivial to prove that
\eqref{acyc:z} and \eqref{acyc:x-y} are equivalent relationships for functions
linked by the valley-index relations. We should note here that, for $\kappa = T$
states, \eqref{acyc:z} and \eqref{acyc:x-y} are not equivalent relationships, so
while $\Phi_{\upsilon x}$ and $\Phi_{\upsilon y}$ are linked they are not linked
with $\Phi_{\upsilon z}$.

When constructing our envelope functions we find it convenient to use variations
of the same underlying naive envelope function for all symmetric states. As
such, we require our naive envelope states to satisfy \eqref{acyc:z}. However,
we note that two such naive envelopes may independently satisfy \eqref{acyc:z}
allowing us freedom to construct $\kappa = T$ states where $\Phi_{\upsilon z}$
is independent of $\Phi_{\upsilon x}$ and $\Phi_{\upsilon y}$. In short, for
$\kappa = T$ states we may have $\Phi_{\upsilon z}$ along with $\Phi_{\upsilon
x} = \left(yzx\right)\Phi_{\upsilon z}^{\prime}$ and $\Phi_{\upsilon y} =
\left(zxy\right)\Phi_{\upsilon z}^{\prime}$ where $\Phi_{\upsilon z} \ne
\Phi_{\upsilon z}^{\prime}$ and both are such that
\begin{align}
\sigma^{\left(\kappa\right)}\left(yxz\right)\Phi_{\upsilon z}
  &=
    \Phi_{\upsilon z}
,\\
\sigma^{\left(\kappa\right)}\left(yxz\right)\Phi_{\upsilon z}^{\prime}
  &=
    \Phi_{\upsilon z}^{\prime}
.
\end{align}

Now we should apply the double inversion factor operator $\mathcal{D}_{z}$ to
the wavefunction. To ensure symmetrization we wish to allow
\begin{align}
  \mathcal{D}_{z}^{\left(\kappa\right)}\psi_{q}^{\prime}
  &=
  \psi_{q}^{\prime}
  .
\label{dblinv}
\end{align}
For \eqref{dblinv} to be true for $q = z$, it must be the case that
\begin{equation}
  \Phi_{\upsilon z}
  =
  \left(\frac{\left(xyz\right)+\left(\bar{x}\bar{y}z\right)}{4}\right)
  \Phi_{\upsilon z}
  +
  \omega^{\left(\kappa\right)}
  \left(
    \frac
    { \left(\bar{x}y\bar{z}\right)+\left(x\bar{y}\bar{z}\right) }
    {4}
  \right)
  \Phi_{-\upsilon z}
  .
\label{dblinv:env:z}
\end{equation}
Some further arithmetic work on \eqref{dblinv:env:z} can establish the following
relationships. We can show that
\begin{align}
  \Phi_{\upsilon z}
  & =
  \left(\bar{x}\bar{y}z\right)
  \Phi_{\upsilon z}
  ,
\label{dblinv:z:para}
\\
  \Phi_{-\upsilon z}
  & =
  \omega^{\left(\kappa\right)}
  \left(
    \frac{\left(\bar{x}y\bar{z}\right)+\left(x\bar{y}\bar{z}\right)}{2}
  \right)
  \Phi_{\upsilon z}
  .
\label{dblinv:z:anti}
\end{align}

We saw in \eqref{acyc:x-y} that the presence of $x$-valley terms implies the
presence of $y$-valley terms in our partially symmetrized wavefunction, and vice
versa. Thus, for the valleys perpendicular to $z$ we must have both
\begin{equation}
  \Phi_{\upsilon x}
  =
  \left(
    \frac
    {
      \left(xyz\right)
      +
      \omega^{\left(\kappa\right)}\left(x\bar{y}\bar{z}\right)
    }
    {4}
  \right)
  \Phi_{\upsilon x}
  +
  \left(
    \frac
    {
      \left(\bar{x}\bar{y}z\right)
      +
      \omega^{\left(\kappa\right)}\left(\bar{x}y\bar{z}\right)
    }
    {4}
  \right)
  \Phi_{-\upsilon x}
  ,
\label{dblinv:x}
\end{equation}
and
\begin{equation}
  \Phi_{\upsilon y}
  =
  \left(
    \frac
    {
      \left(xyz\right)
      +
      \omega^{\left(\kappa\right)}\left(\bar{x}y\bar{z}\right)
    }
    {4}
  \right)
  \Phi_{\upsilon y}
  +
  \left(
    \frac
    {
      \left(\bar{x}\bar{y}z\right)
      +
      \omega^{\left(\kappa\right)}\left(x\bar{y}\bar{z}\right)
    }
    {4}
  \right)
  \Phi_{-\upsilon y}
  .
\label{dblinv:y}
\end{equation}
As before we can do some arithmetic work to simplify these expressions. If we
assume that \eqref{dblinv:x} and \eqref{dblinv:y} are both true we find
\begin{align}
  \Phi_{\upsilon x}
  & =
  \omega^{\left(\kappa\right)}
  \left(x\bar{y}\bar{z}\right)
  \Phi_{\upsilon x}
  ,
\label{dblinv:x:para}
\\
  \Phi_{-\upsilon x}
  & =
  \left(
    \frac
    {
      \left(\bar{x}\bar{y}z\right)
      +
      \omega^{\left(\kappa\right)}\left(\bar{x}y\bar{z}\right)
    }
    {2}
  \right)
  \Phi_{\upsilon x}
  ,
\label{dblinv:x:anti}
\end{align}
as well as
\begin{align}
  \Phi_{\upsilon y}
  & =
  \omega^{\left(\kappa\right)}
  \left(\bar{x}y\bar{z}\right)
  \Phi_{\upsilon y}
  ,
\label{dblinv:y:para}
\\
  \Phi_{-\upsilon y}
  & =
  \left(
    \frac
    {
      \left(\bar{x}\bar{y}z\right)
      +
      \omega^{\left(\kappa\right)}\left(x\bar{y}\bar{z}\right)
    }
    {2}
  \right)
  \Phi_{\upsilon y}
  .
\label{dblinv:y:anti}
\end{align}
It is worth noting that if we assume \eqref{acyc:x-y} is true, and if we also
observe that
\begin{equation}
\left(yxz\right)\left(\bar{x}y\bar{z}\right)\left(yxz\right)
  =
    \left(x\bar{y}\bar{z}\right)
.
\end{equation}
then we see that \eqref{dblinv:x:para} is equivalent to \eqref{dblinv:y:para}
and \eqref{dblinv:x:anti} is equivalent to \eqref{dblinv:y:anti}.

As before, for the $\kappa = A,E$ states we can use the indexing relationships
to try to find additional equivalencies between the above relationships. We once
again invoke \eqref{idx:env:xz}, \eqref{idx:env:yz}, and \eqref{idx:env:yx}. In
that spirit, note the following
\begin{align}
  \left(zxy\right)\left(x\bar{y}\bar{z}\right)\left(yzx\right)
&=
  \left(\bar{x}\bar{y}z\right)
\\
\left(yzx\right)\left(\bar{x}y\bar{z}\right)\left(zxy\right)
&=
  \left(\bar{x}\bar{y}z\right)
\end{align}
Applying these relationships to \eqref{dblinv:x:para} and \eqref{dblinv:y:para}
yields
\begin{equation}
  \Phi_{\upsilon z}
  =
  \omega^{\left(\kappa\right)}
  \left(\bar{x}\bar{y}z\right)
  \Phi_{\upsilon z}
  .
\label{dblinv:xy:para}
\end{equation}

For both \eqref{dblinv:z:para} and \eqref{dblinv:xy:para} to be true, it must be
the case that
\begin{equation}
  \omega^{\left(\kappa\right)} = +1.  \label{AE_cond}
\end{equation}
This condition is satisfied for states with $\kappa = A,E$.

We can also examine the relationships between \eqref{dblinv:z:anti},
\eqref{dblinv:x:anti}, and \eqref{dblinv:y:anti}. Examining the symmetry
operators and their products, we assert that
\begin{align}
  \left(zxy\right)\left(\bar{x}\bar{y}z\right)\left(yzx\right)
  & = 
  \left(\bar{x}y\bar{z}\right)
\\
  \left(zxy\right)\left(\bar{x}y\bar{z}\right)\left(yzx\right)
  & = 
  \left(x\bar{y}\bar{z}\right)
\\
  \left(yzx\right)\left(\bar{x}\bar{y}z\right)\left(zxy\right)
  & = 
  \left(x\bar{y}\bar{z}\right)
\\
  \left(yzx\right)\left(x\bar{y}\bar{z}\right)\left(zxy\right)
  & = 
  \left(\bar{x}y\bar{z}\right)
\end{align}
Applying these relationships to \eqref{dblinv:x:anti}, and \eqref{dblinv:y:anti}
we find
\begin{align}
  \Phi_{-\upsilon z}
  & =
  \left(
    \frac
    {
      \left(\bar{x}y\bar{z}\right)
      +
      \omega^{\left(\kappa\right)}
      \left(x\bar{y}\bar{z}\right)
    }
    {2}
  \right)
  \Phi_{\upsilon z}
  ,
\tag{\ref{dblinv:x:anti}a}
\label{dblinv:x:anti-a}
\\
  \Phi_{-\upsilon z}
  & =
  \left(
    \frac
    {
      \left(x\bar{y}\bar{z}\right)
      +
      \omega^{\left(\kappa\right)}
      \left(\bar{x}y\bar{z}\right)
    }
    {2}
  \right)
  \Phi_{\upsilon z}
  ,
\tag{\ref{dblinv:y:anti}a}
\label{dblinv:y:anti-a}
\end{align}
We note that if \eqref{AE_cond} is true, then \eqref{dblinv:z:anti},
\eqref{dblinv:x:anti-a}, and \eqref{dblinv:y:anti-a} are equivalent. 

So far we have found the form of the partially symmetrized wavefunction to
be such that 
\begin{equation}
  \mathcal{D}_{z}^{\left(\kappa\right)}
  \mathcal{A}_{z}^{\left(\kappa\right)}
  \mathcal{V}_{\kappa,i}^{\dagger}\Psi
  =
  \mathcal{V}_{\kappa,i}^{\dagger}\Psi
  .
\tag{\ref{wf_raw}f}
\label{multi-aleph:form1}
\end{equation}
In essence, the process by which we have applied the $\mathcal{A}_{z}$ acyclic
operator and the $\mathcal{D}_{z}$ double inversion operator hasn't altered the
overall structure of our wavefunction, but has exerted conditions upon the
underlying envelope functions.

It can be shown that the application of the final factor operator $\mathcal{V}$
returns the wavefunction to its initial form. To ensure equality we find the
following requirements upon our envelope functions. For the $\kappa = A,E$
states the following requirements are sufficient.
\begin{align}
  \Phi_{\upsilon z}
  &=
  \sigma^{\left(\kappa\right)}\left(yxz\right)\Phi_{\upsilon z}
  .
\tag{\ref{acyc:z}}
\\
  \Phi_{\upsilon z}
  & =
  \left(\bar{x}\bar{y}z\right)
  \Phi_{\upsilon z}
  ,
\tag{\ref{dblinv:z:para}}
\\
  \Phi_{-\upsilon z}
  & =
  \left(
    \frac{\left(\bar{x}y\bar{z}\right)+\left(x\bar{y}\bar{z}\right)}{2}
  \right)
  \Phi_{\upsilon z}
  .
\tag{\ref{dblinv:z:anti}}
\end{align}
Due to the valley mixing of the $A$- and $E$-states the envelope functions of
other valleys may be found from the following indexing relationships.
\begin{align}
  \Phi_{\upsilon x}
  &=
  3\alpha_{1}^{\left(\kappa,i\right)}\left(yzx\right)
  \Phi_{\upsilon z}
  ,
\tag{\ref{idx:env:xz}}
\\
  \Phi_{\upsilon y}
  &=
  3\alpha_{2}^{\left(\kappa,i\right)}\left(zxy\right)
  \Phi_{\upsilon z}
  .
\tag{\ref{idx:env:yz}}
\end{align}

For the $\kappa = T$ states, the symmetry picks out a longitudinal axis, with
the other axes being perpendicular. The longitudinal envelope functions are
independent of the transverse envelope functions. We can state the requirements
on these envelope functions as such. For the longitudinal envelopes we see
\begin{align}
  \Phi_{\upsilon z}
  &=
  \sigma^{\left(\kappa\right)}\left(yxz\right)\Phi_{\upsilon z}
  .
\tag{\ref{acyc:z}}
\\
  \Phi_{\upsilon z}
  & =
  \left(\bar{x}\bar{y}z\right)
  \Phi_{\upsilon z}
  ,
\tag{\ref{dblinv:z:para}}
\\
  \Phi_{-\upsilon z}
  & =
  -
  \left(
    \frac{\left(\bar{x}y\bar{z}\right)+\left(x\bar{y}\bar{z}\right)}{2}
  \right)
  \Phi_{\upsilon z}
  .
\tag{\ref{dblinv:z:anti}}
\end{align}
The transverse envelopes are independent of the longitudinal envelopes, but not
independent of each other. The relationship between the transverse envelopes
follows 
\begin{equation}
  \Phi_{\upsilon x}
  =
  \sigma^{\left(\kappa\right)}\left(yxz\right)\Phi_{\upsilon y}.
\tag{\ref{acyc:x-y}}
\end{equation}
The transverse envelopes must also follow
\begin{align}
  \Phi_{\upsilon x}
  & =
  \omega^{\left(\kappa\right)}
  \left(x\bar{y}\bar{z}\right)
  \Phi_{\upsilon x}
  ,
\tag{\ref{dblinv:x:para}}
\\
  \Phi_{-\upsilon x}
  & =
  \left(
    \frac
    {
      \left(\bar{x}\bar{y}z\right)
      -
      \left(\bar{x}y\bar{z}\right)
    }
    {2}
  \right)
  \Phi_{\upsilon x}
  .
\tag{\ref{dblinv:x:anti}}
\end{align}